% Autor: Carles Marín (Claude, Anthropic, como asistente de investigación IA).
% Parte II. Qué representaciones de SU(4) contienen la celda de quarks del Modelo Estándar
% en T^2/Z_2: un criterio exacto desde una ramificación de Levi y la proyección quiral del orbifold.
% Narrativa hilada por preguntas; honesta sobre el alcance. Compilar: pdflatex x2.
\documentclass[11pt,a4paper]{article}
\usepackage[utf8]{inputenc}
\usepackage[T1]{fontenc}
\usepackage[spanish,es-noshorthands]{babel}
\usepackage{amsmath,amssymb,amsthm}
\usepackage{booktabs}
\usepackage[table,dvipsnames]{xcolor}
\usepackage{graphicx}
\usepackage{tikz}
\usetikzlibrary{shapes.geometric,arrows.meta,positioning,fit,backgrounds}
\usepackage{geometry}
\geometry{margin=2.4cm}
\usepackage[colorlinks=true,linkcolor=RoyalBlue,citecolor=BrickRed,urlcolor=RoyalBlue]{hyperref}
\usepackage{caption}
\captionsetup{font=small,labelfont=bf}
\newtheorem{teorema}{Teorema}
\newtheorem{lema}{Lema}
\newcommand{\ZZ}{\mathbb{Z}}
\newcommand{\RR}{\mathbb{R}}
\definecolor{okgreen}{RGB}{0,140,54}
\definecolor{hdrblue}{RGB}{31,78,121}
\definecolor{softgrey}{RGB}{238,242,247}

\title{\textbf{Tres Puertas a una Generación de Quarks}\\[3pt]
\textbf{Un Criterio Exacto para Qué Representaciones de $SU(4)$ Contienen el Modelo Estándar,\\
y Por Qué el Grupo Gauge Es Crítico en Rango Tres}\\[5pt]
\large\textnormal{Parte II --- la celda $\pm\tfrac12$, una ramificación de Levi, y la
proyección quiral del orbifold}}
\author{Carles Mar\'in\thanks{Investigador independiente, \texttt{karlesmarin@gmail.com}.
Los cálculos se realizaron y contrastaron con Claude (Anthropic) como asistente de
investigación IA contra una verdad-base verificable por máquina (aritmética racional exacta
sobre los sistemas de pesos de $SU(n)$ en SageMath); todo recuento de este trabajo se
regenera desde los scripts anexos.}}
\date{\today\\[2pt]\small DOI: \href{https://doi.org/10.5281/zenodo.21432628}{10.5281/zenodo.21432628}}

\begin{document}
\maketitle

\begin{abstract}
\noindent
El trabajo compañero embebió una generación completa de quarks del Modelo Estándar en la
representación de dimensión $60$, $(3,\mathbf{60})$, de $SU(4)$ sobre $T^2/\mathbb{Z}_2$ y
mostró, por un barrido exhaustivo, que $(3,\mathbf{60})$ es la representación \emph{más
pequeña} capaz de hacerlo. Esa minimalidad invita a la pregunta general: \emph{cuáles}
representaciones de $SU(4)$ pueden alojar una generación de quarks, y por qué solo esas. La
respondemos exactamente. Una representación irreducible de labels de Dynkin $(a,b,c)$
contiene la celda de quarks del Modelo Estándar
$\{Q(\mathbf{2},\tfrac16),\,u(\mathbf{1},\tfrac23),\,d(\mathbf{1},-\tfrac13)\}$ en sus modos
cero quirales de $T^2/\mathbb{Z}_2$ si y solo si
\[
(a+2b+3c)\ \text{impar}\quad\wedge\quad b\ge 1\quad\wedge\quad a+b+c\ge 3 .
\]
Las tres cláusulas son, en orden, una condición \emph{aritmética} (la carga del centro
$\ZZ_4$), una \emph{geométrica} (el nodo central de $SU(4)$ excitado) y una de \emph{tamaño}
(sitio suficiente para cerrar la celda). Derivamos el criterio borrando el nodo central del
diagrama de Dynkin de $SU(4)$ --- exponiendo un subgrupo de Levi
$SU(2)_L\times SU(2)_R\times U(1)$ --- y proyectando la torre de Clebsch--Gordan resultante
sobre los modos cero quirales del orbifold; la proyección colapsa la torre a un recuento
cerrado de modos cero $N=(b+1)(a+c+1)/2$. Somos explícitos sobre qué es nuevo y qué es
heredado: la condición del centro es clásica, la ramificación es de Littlewood--Richardson, y
la propiedad del punto medio (el $\tfrac16$ del doblete es la media de $\tfrac23$ y $-\tfrac13$
de los singletes) refleja la relación de hipercarga de Pati--Salam; lo nuestro es la
proyección quiral que produce el recuento cerrado, el empaquetado del criterio, y la lectura
de la celda como el objeto invariante. Finalmente situamos a $SU(4)$ en una familia: la celda
impone tres restricciones de carga sobre un funcional de hipercarga con $\mathrm{rango}(G)-1$
parámetros libres, por lo que es una obstrucción \emph{rígida} exactamente en
$\mathrm{rango}=3$ --- $SU(4)$ es el caso marginal, y para rango mayor una representación
genérica admite. El mismo recuento explica, como un efecto de grados-de-libertad-consumidos,
por qué una hipercarga fijada en modelos de rango mayor deja a los quarks con las cargas
(enteras) equivocadas y los fuerza a una brana.
\end{abstract}

% ===================================================================
\section{De una representación a todas ellas}\label{sec:intro}
Este trabajo continúa una historia que su compañero comenzó. La Parte~I~\cite{PartI} retomó
una pregunta abierta desde 2003, cuando Scrucca, Serone, Silvestrini y Wulzer observaron que
en un orbifold la anulación de los tadpoles localizados y la cancelación de las anomalías son
requisitos \emph{independientes}, y nombraron la pieza que faltaba --- un espectro libre de
anomalías con tadpole globalmente nulo --- como un problema abierto~\cite{SSSW}. La Parte~I lo
abordó en el modelo six-dimensional de dos dobletes de Higgs con $SU(4)$ de Akamatsu, Hirose,
Maru y Nago~\cite{AHMN}: embebió un bloque coloreado de quarks en la representación de
dimensión $60$, $(3,\mathbf{60})$, exhibió una completación de $45$ multipletes que cancela
anomalías y tadpole juntos, y demostró un \emph{Teorema de Existencia} --- sobre esa librería
fermiónica el tadpole nunca obstruye una completación libre de anomalías --- convirtiendo un
testigo computacional en una garantía estructural. Cerró con una propuesta: allí donde un
resultado de construcción de modelos diga ``hemos encontrado uno'', preguntar si el hallazgo
era en realidad una ley. Este trabajo toma a la Parte~I por su palabra. Pues de pasada, y sin
explicar por qué, la Parte~I había registrado un segundo hallazgo --- un barrido exhaustivo
mostrando que $(3,\mathbf{60})$ es la representación \emph{más pequeña} de $SU(4)$ capaz de
alojar una generación completa de quarks. ¿Por qué esa, y por qué tan pocas? Convertimos el
hecho barrido en una ley exacta.

Esa pregunta es más antigua que el modelo, y su historia es la de cajas empaquetadas cada vez
más apretadas. Una construcción gran-unificada o de gauge--Higgs comienza eligiendo una
representación que porte la materia, y la elección está severamente restringida: la
representación debe contener, entre sus estados, exactamente los números cuánticos de una
generación. Georgi y Glashow encajaron una generación en el
$\overline{\mathbf 5}\oplus\mathbf{10}$ de $SU(5)$~\cite{GeorgiGlashow} --- en dos piezas, y
sin asiento para un neutrino de mano derecha. Un año después el espinor $\mathbf{16}$ de
$SO(10)$ cerró la caja: quince fermiones del Modelo Estándar, y un decimosexto lugar que el
espinor simplemente \emph{tiene}, casi obligando a sentar allí un $\nu_R$~\cite{SO10FM,SO10G}
--- el grupo sabiendo, por así decir, antes que nosotros. El excepcional $\mathbf{27}$ de
$E_6$ siguió, ofreciendo una generación más exóticos~\cite{E6}, y Georgi pidió abiertamente
las tres familias en una sola representación irreducible~\cite{GeorgiFlavor,Yamatsu}, un sueño
que naufraga en un salón de fermiones espejo. A lo largo de todo ello las tablas de Slansky
fueron la guía telefónica del campo para decidir qué contiene una
representación~\cite{Slansky}; el censo reciente de Herms y Ruhdorfer convierte los recuentos
en una estadística --- qué fracción de la materia libre de anomalías es
unificable~\cite{HermsRuhdorfer}. Pero la decisión siguió siendo, representación a
representación, una consulta y no una ley.

Un mínimo es una invitación. Si una representación es la más pequeña, hay un paisaje detrás de
ella --- y la Parte~I cartografió un único punto. ¿Qué distingue a las representaciones que
tienen éxito de las que fracasan? ¿Es ``$(3,\mathbf{60})$'' un accidente de dimensión, o el
primer miembro de una familia con una regla limpia? Este trabajo da la regla, exactamente, y
lee de ella por qué la familia es tan escasa, por qué el mínimo cae donde cae, y por qué el
grupo gauge es $SU(4)$ y no algo mayor. Donde la Parte~I validó un testigo, la Parte~II
cartografía el espacio a su alrededor; juntas van de un modelo consistente a la comprensión de
todo el paisaje de materia en que vive.

El objeto organizador en todo momento es lo que llamaremos la \emph{celda del Modelo
Estándar}: el contenido de quarks de triplete de color de una generación, reducido a su
esqueleto electrodébil,
\[
Q\,(\mathbf{2},\tfrac16),\qquad u\,(\mathbf{1},\tfrac23),\qquad d\,(\mathbf{1},-\tfrac13).
\]
Los singletes up y down son el doblete desplazado en hipercarga por $\pm\tfrac12$ --- el
isospín de mano derecha $T^3_R$ de la relación de Pati--Salam que encontramos en
\S\ref{sec:cell}. El hecho operativo es una pequeña pieza de aritmética: el $\tfrac16$ del
doblete es la \emph{media} de las cargas de los singletes,
$\tfrac16=\tfrac12\!\big(\tfrac23+(-\tfrac13)\big)$, de modo que el doblete se sitúa en el
punto medio exacto de hipercarga de sus dos singletes. Ese punto medio es lo que hace de la
celda un objeto geométrico rígido en el retículo de pesos, y es donde comienza nuestro
criterio.

\paragraph{Alcance y relación con trabajo previo.} Somos cuidadosos en separar lo nuevo de lo
heredado, y la mayoría de los ingredientes son heredados. La condición de carga del centro ---
que admitir las hipercargas fraccionarias del Modelo Estándar fuerza una clase de congruencia
fija del centro del grupo gauge --- es clásica, enunciada para el propio $\ZZ_6$ del Modelo
Estándar por Hucks y por Saller y revisitada por Tong y por Alonso, Dimakou y
West~\cite{Hucks,Saller,Tong,Alonso}. La ramificación que usamos es una restricción de Levi
parabólica maximal, cuyas multiplicidades son coeficientes de
Littlewood--Richardson~\cite{Fulton,GW}. La relación del punto medio es la identidad de
hipercarga de Pati--Salam, y la observación de que un funcional de hipercarga embebido en un
grupo simple tiene un número fijo de parámetros libres es el mismo recuento que subyace a la
clásica cota inferior ``$\mathrm{rango}\ge 4$'' sobre grupos unificados; un análisis
estrechamente relacionado de funcional-de-Cartan de modelos gauge-Higgs, aunque del sector de
bosones gauge y no de la celda de materia, se debe a Aranda y Wudka~\cite{ArandaWudka}. No
reclamamos nueva teoría de grupos. Lo nuestro es triple: (i) el \emph{criterio exacto de tres
condiciones} para $SU(4)$ sobre $T^2/\mathbb{Z}_2$, empaquetado como un enunciado cerrado
único (\S\ref{sec:theorem}); (ii) la \emph{proyección quiral del orbifold} que colapsa la
torre de Clebsch--Gordan de Levi al recuento cerrado de modos cero $N=(b+1)(a+c+1)/2$
(\S\ref{sec:levi}); y (iii) la lectura de la celda $\pm\tfrac12$ como el invariante, junto con
la ubicación de $SU(4)$ en el rango crítico (\S\ref{sec:rank}) --- lo último una lente
interpretativa sobre ingredientes clásicos, no un teorema nuevo. La Figura~\ref{fig:roadmap}
es la ruta.

\begin{figure}[t]
\centering
\resizebox{\textwidth}{!}{%
\begin{tikzpicture}[font=\small,
  act/.style={rounded corners,draw,thick,text width=3.35cm,align=center,inner sep=6pt,minimum height=4.0cm,anchor=north},
  arr/.style={-{Stealth[length=3.4mm]},very thick,black!55}]
\node[act,fill=RoyalBlue!8,draw=RoyalBlue!65] (a1) at (0,0)
 {\textbf{I.\ Una}\\[1pt]{\footnotesize\S\,1}\\[5pt]
  \textit{$(3,\mathbf{60})$ aloja una}\\\textit{generación, y es}\\\textit{mínima (Parte I)}\\[6pt]
  {\footnotesize\color{RoyalBlue}\textbf{un solo testigo}}};
\node[act,fill=okgreen!9,draw=okgreen!65] (a2) at (5.5,0)
 {\textbf{II.\ Todas}\\[1pt]{\footnotesize\S\S\,2--5}\\[5pt]
  \textit{¿Qué $(a,b,c)$}\\\textit{contienen la celda?}\\[6pt]
  {\footnotesize$(a{+}2b{+}3c)$ impar\\$\wedge\ b\ge1\ \wedge\ a{+}b{+}c\ge3$}\\[5pt]
  {\footnotesize\color{okgreen}\textbf{criterio exacto}}};
\node[act,fill=softgrey,draw=hdrblue!65] (a3) at (11.0,0)
 {\textbf{III.\ Por qué}\\[1pt]{\footnotesize\S\S\,4,6}\\[5pt]
  \textit{Borra el nodo}\\\textit{central: una torre de}\\\textit{Levi, partida por paridad}\\[6pt]
  {\footnotesize $N=(b{+}1)\tfrac{a+c+1}{2}$}};
\node[act,fill=black!5,draw=black!45,dashed] (a4) at (16.5,0)
 {\textbf{IV.\ Qué grupo}\\[1pt]{\footnotesize\S\,7}\\[5pt]
  \textit{Rígido solo en}\\\textit{$\mathrm{rango}=3$;}\\\textit{$SU(4)$ marginal}\\[6pt]
  {\footnotesize\color{hdrblue}\textbf{caso crítico}}};
\draw[arr] (a1.east) -- node[above,font=\footnotesize\itshape,black!70]{generaliza} (a2.west);
\draw[arr] (a2.east) -- node[above,font=\footnotesize\itshape,black!70]{deriva} (a3.west);
\draw[arr] (a3.east) -- node[above,font=\footnotesize\itshape,black!70]{ubica} (a4.west);
\end{tikzpicture}}
\caption{La ruta de este trabajo. Del único testigo mínimo de la Parte~I a un criterio exacto
para todas las representaciones; el criterio se deriva por una ramificación de Levi y una
proyección quiral, y $SU(4)$ queda ubicado como el único rango marginal en el que la celda es
una obstrucción rígida.}
\label{fig:roadmap}
\end{figure}

% ===================================================================
\section{La celda del Modelo Estándar como un problema de retículo}\label{sec:cell}
¿Qué significa, exactamente, que una representación \emph{contenga} el Modelo Estándar --- y
podemos decidirlo por un cálculo en vez de por una búsqueda entre estados? Para afilar la
pregunta primero reducimos una generación al objeto que debe estar presente. El color no
juega papel activo en lo que sigue: en el modelo de la Parte~I el $SU(3)_c$ fuerte es un
factor espectador y el bloque de quarks es $(\mathbf 3, R)$ con $R$ una representación de
$SU(4)$, de modo que ``triplete de color'' es automático y el contenido electrodébil vive
enteramente en $R$. El $SU(2)_L$ débil está generado por la primera raíz simple de $SU(4)$, y
la hipercarga es una dirección de Cartan que conmuta con él. Escribiendo un peso de $SU(4)$ en
las coordenadas de partición $(n_1,n_2,n_3,n_4)$ y las dos cargas $SU(2)_L$-invariantes
independientes
\[
q_8=n_1+n_2-2n_3,\qquad q_{15}=n_1+n_2+n_3-3n_4,
\]
la hipercarga no rota que reproduce las cargas de quarks del Modelo Estándar es
$Y=(-7q_8+4q_{15})/18$~\cite{PartI}, un elemento específico de la familia de dos parámetros de
funcionales $SU(2)_L$-invariantes sin traza $Y=\alpha q_8+\beta q_{15}$.

Una representación \emph{contiene la celda} si sus modos cero incluyen un doblete débil de mano
izquierda $Q$ y dos singletes débiles de mano derecha $u,d$ cuyas cargas bajo un único tal $Y$
son $\tfrac16,\tfrac23,-\tfrac13$. Como $Y$ es lineal y homogénea, esto es una condición exacta
de álgebra lineal. Recogiendo los vectores de carga de los tres multipletes, la celda cierra
precisamente cuando
\begin{equation}\label{eq:det}
\det\!\begin{pmatrix} q_8^{Q} & q_{15}^{Q} & \tfrac16\\[2pt]
q_8^{u} & q_{15}^{u} & \tfrac23\\[2pt] q_8^{d} & q_{15}^{d} & -\tfrac13\end{pmatrix}=0
\qquad\text{con el menor $2\times2$ superior no nulo,}
\end{equation}
es decir, el vector objetivo $(\tfrac16,\tfrac23,-\tfrac13)$ yace en el subespacio generado por
las dos columnas de carga, de modo que una única $Y$ fija los tres valores. La
Ecuación~\eqref{eq:det} convierte ``contiene el Modelo Estándar'' en un enunciado sobre la
geometría afín del diagrama de pesos.

Una consecuencia es inmediata y organizadora. Como
$\tfrac16=\tfrac12\big(\tfrac23+(-\tfrac13)\big)$, la carga del doblete es la \emph{media} de
las dos cargas de los singletes; si el doblete se sitúa en el punto medio $Q=\tfrac12(u+d)$ de
los singletes en el espacio de cargas, entonces $Y(Q)=\tfrac16$ se cumple
\emph{automáticamente} una vez $Y(u)=\tfrac23$ y $Y(d)=-\tfrac13$. La celda es un par de
reflexión simétrico $\pm\tfrac12$ a ambos lados del doblete --- la cara geométrica de la
relación de hipercarga de Pati--Salam, de medio siglo de antigüedad, $Y=T^3_R+(B-L)/2$
\cite{PatiSalam} (con $Q_{\rm em}=T^3+Y$, de modo que $T^3_R=0$ para el doblete y $\pm\tfrac12$
para los singletes, y $B-L=\tfrac13$ para un quark); el mismísimo $\pm\tfrac12$ que da nombre a
la celda es el isospín de mano derecha de Pati y Salam.\footnote{Una advertencia contra un
juego de palabras natural: el grupo que porta nuestra celda es el $SU(4)$ electrodébil del
modelo AHMN, con el color un $SU(3)_c$ \emph{separado}; no es el $SU(4)_c$ de color de Pati y
Salam, en el que el número leptónico es un cuarto color. Lo que tomamos prestado
de~\cite{PatiSalam} es la relación de hipercarga, no el grupo.} El problema es ahora: ¿para qué
$(a,b,c)$ ofrece el diagrama de pesos de $R$ tal configuración entre sus modos cero
seleccionados por paridad?

\emph{Establecido aquí:} contener el Modelo Estándar se reduce a una condición exacta de
rango~\eqref{eq:det} sobre el diagrama de pesos, con la celda $\pm\tfrac12$ un par de reflexión
simétrico. \emph{Abierto:} qué labels $(a,b,c)$ ofrecen de hecho esa configuración --- las tres
condiciones del resto del trabajo.

% ===================================================================
\section{La condición del centro (clásica)}\label{sec:centre}
Antes de preguntar \emph{dónde} se sitúan las cargas, una pregunta más contundente: ¿qué
representaciones pueden mostrar las fracciones $\tfrac16,\tfrac23,\tfrac13$ del Modelo Estándar
\emph{siquiera}? La primera obstrucción es aritmética y no es nuestra. El valor de una
hipercarga sobre un peso es, salvo normalización, una combinación entera de los labels dividida
por el orden del centro; las fracciones $\tfrac16,\tfrac23,\tfrac13$ del Modelo Estándar solo
pueden alcanzarse desde una clase de congruencia fija. Para $SU(4)$ el invariante relevante es
la carga del centro $\ZZ_4$, igual a la n-alidad
\begin{equation}
c(R)=a+2b+3c \pmod 4,
\end{equation}
la clase de congruencia de Slansky~\cite{Slansky}. Una comprobación directa sobre la hipercarga
anterior muestra que las fracciones de los quarks ocurren solo cuando esta clase es
\emph{impar}. Que la hipercarga fraccionaria esté atada a la carga del centro exactamente de
este modo es un hecho clásico, enunciado para el propio $\ZZ_6=\ZZ_2\times\ZZ_3$ del Modelo
Estándar por Hucks y por Saller, y reformulado en lenguaje de simetría de 1-forma /
estructura-global por Tong y por Alonso, Dimakou y West~\cite{Hucks,Saller,Tong,Alonso}. No
reclamamos nada de ello; registramos que, transferido al centro $\ZZ_4$ del grupo de embedding,
es la primera de nuestras tres cláusulas.

\emph{Establecido aquí:} la condición uno, $(a+2b+3c)$ impar, es la congruencia clásica de
carga del centro, heredada no inventada. \emph{Abierto:} las dos condiciones geométricas, que
el diagrama de pesos --- no el centro --- debe suministrar.

% ===================================================================
\section{Borra el nodo central: la torre de Levi y su proyección quiral}\label{sec:levi}
El centro dice que las fracciones son alcanzables. Pero, ¿alcanzables \emph{juntas}, bajo una
única hipercarga --- y dónde, en $SU(4)$ mismo, se decide eso? Las dos cláusulas restantes, y
la razón de ellas, provienen de un único movimiento estructural. El diagrama de Dynkin de
$SU(4)$ es un camino de tres nodos, $\circ\!-\!\circ\!-\!\circ$; \emph{borrar el nodo central}
lo parte en dos $SU(2)$ disjuntos y deja un $U(1)$:
\begin{equation}
SU(4)\ \supset\ \underbrace{SU(2)_L}_{\text{nodo }\alpha_1}\ \times\
\underbrace{SU(2)_R}_{\text{nodo }\alpha_3}\ \times\ \underbrace{U(1)_m}_{\text{borrado }\alpha_2}.
\end{equation}
$SU(2)_L$ es el isospín débil. El nodo borrado porta la carga $U(1)$
\begin{equation}
m\ :=\ 2q_8+q_{15}\ =\ 3(n_1+n_2-n_3-n_4),
\end{equation}
que --- como se comprueba sobre las raíces simples --- solo cambia bajo la raíz central
$\alpha_2$ y queda fija bajo $\alpha_1,\alpha_3$: es la coordenada conjugada al nodo central.
Esta única observación controla las dos condiciones restantes.

Ramifica una irreducible $(a,b,c)$ a este Levi. Su sector de \emph{singletes} de $SU(2)_L$ ---
el asiento de los singletes de quarks de mano derecha $u,d$ --- tiene una forma limpia
(Figura~\ref{fig:levi}, verificada en todas las representaciones hasta dimensión $700$):
consiste en
\begin{equation}\label{eq:tower}
(b+1)\ \text{copias de la torre de Clebsch--Gordan}\quad
[\tfrac a2]\otimes[\tfrac c2]\ =\ \bigoplus_{j=|a-c|/2}^{(a+c)/2}[j],
\end{equation}
una copia en cada uno de los $b+1$ valores equiespaciados de $m$ (espaciado $12$, de modo que
los singletes abarcan una extensión de nodo central de exactamente $12b$). El número de copias
lo fija el label central $b$; el contenido de espín de $SU(2)_R$ interno lo fijan los labels
externos $a,c$, con espín \emph{máximo} $(a+c)/2$.

Ahora impón el orbifold. Sobre $T^2/\mathbb{Z}_2$ las dos paridades $\ZZ_2$ asignan quiralidad;
en el sector de singletes están bloqueadas (una es redundante), de modo que los modos cero de
mano derecha supervivientes son exactamente la \emph{mitad} de cada torre, y --- como las
hipercargas distintas las fija la \emph{extensión} de la torre, es decir su espín máximo
$(a+c)/2$ --- el recuento de huecos de carga distintos de singlete de color de mano derecha es
\begin{equation}\label{eq:count}
\boxed{\,N(a,b,c)\ =\ (b+1)\cdot\frac{a+c+1}{2}\,}
\end{equation}
(el factor externo es entero porque $a+c$ es impar siempre que la carga del centro es impar).
La Ecuación~\eqref{eq:count} es el corazón técnico del trabajo: un recuento cerrado de
multipletes de modos cero quirales como función de los labels de Dynkin de la representación
bulk. El recuento factoriza como
\[
N=\underbrace{(b+1)(a+c+1)}_{\text{ramificación de Levi (indep.\ de dimensión)}}\ \times\ \underbrace{\tfrac12}_{\text{proyección de orbifold}},
\]
una forma que hace transparentes los papeles: la teoría de grupos es heredada (una evaluación
de Littlewood--Richardson~\cite{Fulton,GW}), mientras que el partir por la mitad es la
proyección extra-dimensional. Es la \emph{proyección quiral extremal} --- quédate con el espín
máximo, parte por paridad --- lo que es nuestro, no la forma cerrada que produce.

\emph{Establecido aquí:} un recuento cerrado de modos cero quirales $N=(b+1)(a+c+1)/2$, con
extensión de nodo central $12b$, desde una ramificación de Levi y la proyección de paridad.
\emph{Abierto:} cómo decide este recuento si la celda cierra --- el criterio de
\S\ref{sec:theorem}.

\begin{figure}[t]
\centering
\begin{tikzpicture}[font=\small,x=1.05cm,y=0.92cm]
\begin{scope}[shift={(1.6,4.25)}]
  \node[circle,draw,thick,fill=RoyalBlue!20,minimum size=6mm] (d1) at (0,0){};
  \node[circle,draw,thick,dashed,fill=black!6,minimum size=6mm] (d2) at (1.1,0){};
  \node[circle,draw,thick,fill=okgreen!25,minimum size=6mm] (d3) at (2.2,0){};
  \draw[thick] (d1)--(d2)--(d3);
  \node[below=1pt of d1,font=\footnotesize]{$\alpha_1$};
  \node[below=1pt of d2,font=\footnotesize]{$\alpha_2$};
  \node[below=1pt of d3,font=\footnotesize]{$\alpha_3$};
  \node[above=1pt of d1,font=\scriptsize,RoyalBlue]{$SU(2)_L$};
  \node[above=8pt of d2,font=\scriptsize,black!60]{$U(1)_m$ (borr.)};
  \node[above=1pt of d3,font=\scriptsize,okgreen!60!black]{$SU(2)_R$};
\end{scope}
\draw[->,black!55,thick] (-0.6,-0.2) -- (7.4,-0.2) node[right,font=\footnotesize]{$q_8$ (transversal)};
\draw[->,black!55,thick] (-0.4,-0.4) -- (-0.4,3.9) node[above,font=\footnotesize]{$m=2q_8+q_{15}$};
\foreach \lev/\yy in {0/0.2, 1/1.2, 2/2.2, 3/3.2}{
  \node[left,font=\scriptsize,black!55] at (-0.5,\yy){$m_\lev$};
}
\foreach \yy/\xa/\xb in {0.2/0.4/1.6, 1.2/1.6/2.8, 2.2/2.9/4.1, 3.2/4.2/5.4}{
  \fill[okgreen!70!black] (\xa,\yy) circle (2.6pt);
  \fill[okgreen!70!black] (\xb,\yy) circle (2.6pt);
  \draw[black!25] (\xa,\yy)--(\xb,\yy);
}
\fill[BrickRed] (2.25,1.7) circle (2.9pt);
\node[BrickRed,font=\scriptsize,above right] at (2.25,1.7){$Q\,(\tfrac16)$};
\node[okgreen!55!black,font=\scriptsize,below] at (1.6,1.2){$u\,(\tfrac23)$};
\node[okgreen!55!black,font=\scriptsize,above] at (2.9,2.2){$d\,(-\tfrac13)$};
\draw[BrickRed,dashed,thick] (1.6,1.2)--(2.9,2.2);
\end{tikzpicture}
\caption{Borrar el nodo central de $SU(4)$ expone $SU(2)_L\times SU(2)_R\times U(1)_m$. Los
modos cero de singlete de color de mano derecha (verde) forman un retículo cizallado: $b+1$
niveles en la carga del nodo central $m$ (aquí $b=3$, cuatro niveles), cada uno portando la
mitad de una torre de Clebsch--Gordan. Una celda del Modelo Estándar (rojo/verde, discontinua)
es un par simétrico $\pm\tfrac12$ $u,d$ cuyo punto medio es un doblete débil $Q$ en hipercarga
$\tfrac16$. La celda solo puede cerrar cuando el retículo es genuinamente bidimensional --- lo
que requiere el nodo central excitado ($b\ge1$) y al menos tres singletes ($a+b+c\ge3$).}
\label{fig:levi}
\end{figure}

% ===================================================================
\section{El criterio de la celda}\label{sec:theorem}
Ya disponemos de un recuento exacto de los singletes; ¿cuándo basta un recuento para cerrar una
celda? El recuento~\eqref{eq:count} decide las dos últimas condiciones. La celda~\eqref{eq:det}
necesita un retículo de singletes \emph{bidimensional}: tres huecos de carga no colineales, de
modo que algún par simétrico tenga un doblete en su punto medio. Leyéndolo de la torre:
\begin{itemize}
\item si $b=0$ hay un único valor de $m$: todos los singletes son colineales, y ninguna
hipercarga puede separar la celda. Por tanto $b\ge1$ es necesario --- el nodo central debe
excitarse para abrir la dirección transversal (espín-$SU(2)_R$); la extensión es exactamente
$12b$.
\item con $b\ge1$, el número de singletes distintos es $N=(b+1)(a+c+1)/2$; esto cae a $N=2$
solo en la única esquina $b=1,\ a+c=1$, i.e.\ $a+b+c=2$, donde los dos singletes disponibles
son colineales y la celda está sobre-restringida. Para todos los demás casos $N\ge3$ y el
retículo es bidimensional. Por tanto $a+b+c\ge3$.
\end{itemize}
Junto con la condición del centro esto da el criterio, que hemos verificado exhaustivamente
contra la construcción directa~\eqref{eq:det} en todas las representaciones hasta dimensión
$900$.

\begin{teorema}[Criterio de la celda]\label{thm:cell}
Una representación irreducible de dimensión finita $(a,b,c)$ de $SU(4)$ contiene una celda
completa de quarks del Modelo Estándar
$\{Q(\mathbf 2,\tfrac16),u(\mathbf 1,\tfrac23),d(\mathbf 1,-\tfrac13)\}$ entre sus modos cero
quirales de mano derecha/izquierda de $T^2/\mathbb{Z}_2$ si y solo si
\[
\boxed{\;(a+2b+3c)\ \text{es impar}\quad\wedge\quad b\ge1\quad\wedge\quad a+b+c\ge3\;}
\]
Equivalentemente: la carga del centro es impar (aritmética), el nodo central está excitado de
modo que el retículo de singletes de color de mano derecha se extiende $12b$ en la carga del
nodo central (geométrica), y el retículo porta al menos tres puntos no colineales
$N=(b+1)(a+c+1)/2\ge3$ (tamaño).
\end{teorema}

\textbf{El Teorema~\ref{thm:cell} es el resultado central de este trabajo, y el criterio es
nuestro}: el enunciado cerrado de tres condiciones, y el recuento por proyección
quiral~\eqref{eq:count} que lo produce, son nuevos. Lo decimos sin sobre-reclamar --- la
Tabla~\ref{tab:gates} registra las tres cláusulas y, junto a cada una, exactamente qué se
hereda de la teoría de grupos clásica y qué es nuevo, de modo que la autoría es explícita
condición a condición. El criterio es escaso: hasta dimensión $400$ solo ocho representaciones
califican (salvo conjugación), de dimensiones $60,84,140,140,216,224,280,360$, y la más
pequeña, $(a,b,c)=(0,2,1)$ de dimensión $60$, es precisamente la $(3,\mathbf{60})$ de la
Parte~I --- de modo que el Teorema~\ref{thm:cell} recupera, y explica, la minimalidad hallada
allí por fuerza bruta.

\begin{table}[t]
\centering\small
\begin{tabular}{p{3.4cm}p{4.9cm}p{5.4cm}}
\toprule
\rowcolor{hdrblue}
\textcolor{white}{Condición} & \textcolor{white}{Heredado (citar)} & \textcolor{white}{Nuestro}\\
\midrule
$(a{+}2b{+}3c)$ impar\newline\emph{(aritmética)} &
congruencia del centro $\ZZ_4$ / n-alidad; $Y$ fraccionaria atada al
centro~\cite{Hucks,Saller,Tong,Alonso,Slansky} &
transferencia al $\ZZ_4$ del grupo de embedding como condición 1 de un criterio cerrado
único\\
\midrule
\rowcolor{softgrey}
$b\ge1$\newline\emph{(geométrica)} &
ramificación de Levi / torre de Clebsch--Gordan~\cite{Fulton,GW} &
la extensión de nodo central $=12b$; su necesidad para un doblete quiral en $Y=\tfrac16$\\
\midrule
$a+b+c\ge3$\newline\emph{(tamaño)} &
--- &
el recuento cerrado de modos cero quirales $N=(b{+}1)(a{+}c{+}1)/2$ y su umbral $N\ge3$\\
\bottomrule
\end{tabular}
\caption{Las tres condiciones del Teorema~\ref{thm:cell}, y un balance honesto de qué debe cada
una a la teoría de grupos clásica y qué es nuevo aquí. La novedad es la proyección quiral que da
el recuento cerrado, y el empaquetado de las tres en un criterio --- no la aritmética del centro
ni la ramificación.}
\label{tab:gates}
\end{table}

\begin{figure}[t]
\centering
\includegraphics[width=0.86\linewidth]{fig_cell_landscape_es.pdf}
\caption{El paisaje de admisibilidad del Teorema~\ref{thm:cell}, sobre el plano de labels
$(a+c,\,b)$ (la condición del centro restringe a $a+c$ impar). El fondo es el recuento exacto
de modos cero $N=(b+1)(a+c+1)/2$; un anillo verde marca una representación que pasa las tres
condiciones, una cruz roja una que falla. Las condiciones recortan la región admisible con
limpieza: cae toda la fila $b=0$ (nodo central no excitado), como cae la única esquina
$a+c=1,\,b=1$ (demasiado pequeña, $N=2$). El mínimo, $(0,2,1)=\mathbf{60}$ de la Parte~I, se
sitúa en la esquina inferior izquierda de la región admisible. Cada marcador está calculado, no
es esquemático.}
\label{fig:landscape}
\end{figure}

\begin{figure}[t]
\centering
\includegraphics[width=0.80\linewidth]{fig_cell_lattice_es.pdf}
\caption{El retículo de pesos de singletes de color de mano derecha de una representación
admisible concreta, $(0,3,3)$, calculado exactamente. Los singletes (relleno, coloreados por la
carga del nodo central $m=2q_8+q_{15}$) forman el retículo cizallado
$(b{+}1)\times\frac{a+c+1}{2}$ de \S\ref{sec:levi} --- aquí cuatro niveles de $m$ ($b=3$). Los
cuadrados grises tenues son los dobletes débiles. Una celda del Modelo Estándar (rojo) es un par
simétrico $\pm\tfrac12$ $u,d$ cuyo punto medio es un doblete $Q$; el cizallamiento es
exactamente lo que permite que tres puntos no sean colineales, para que la celda pueda cerrar.
Estos son los datos tras el esquema de la Figura~\ref{fig:levi}.}
\label{fig:latticereal}
\end{figure}

% ===================================================================
\section{Un inciso sobre la independencia del nodo}\label{sec:nodeindep}
¿Por qué es tan limpio ese recuento --- es $(a{+}1)(b{+}1)(c{+}1)$ un accidente de $SU(4)$, o
insiste el grupo en ello? La forma limpia de~\eqref{eq:count} descansa en un pequeño hecho
general que vale la pena aislar, aunque solo sea para ver dónde acaba la pulcritud de $SU(4)$.
Para una raíz simple $\alpha$, escribe $\dim\mathrm{Inv}_\alpha(V)$ para la dimensión del
subespacio de una representación $V$ aniquilado por el $SU(2)$ generado por $\alpha$.

\begin{lema}\label{lem:node}
Para cualquier representación irreducible de un álgebra de Lie simple, $\dim\mathrm{Inv}_\alpha(V)$
depende solo de la órbita de Weyl --- equivalentemente la longitud --- de la raíz simple
$\alpha$. En un álgebra simply-laced (todas las raíces de una longitud, luego una órbita de
Weyl) es por tanto independiente de qué raíz simple se elija; en un álgebra no-simply-laced toma
a lo sumo dos valores, uno para raíces largas y otro para cortas.
\end{lema}

\begin{proof}
Las multiplicidades de peso son Weyl-invariantes. Para $w$ en el grupo de Weyl,
$\langle w\mu,(w\alpha)^\vee\rangle=\langle\mu,\alpha^\vee\rangle$, de modo que la distribución
de $\alpha$-autovalores sobre el multiconjunto (Weyl-estable) de pesos es la misma para $\alpha$
y $w\alpha$; luego también lo es $\dim\mathrm{Inv}_\alpha=N_\alpha(0)-N_\alpha(2)$. En un álgebra
simply-laced todas las raíces yacen en una órbita de Weyl.
\end{proof}

Para $SU(4)$ (tipo $A_3$, simply-laced) la dimensión de invariantes es así la misma en cada
nodo, e igual al producto limpio $(a+1)(b+1)(c+1)$. Esta forma-producto, sin embargo, es
especial de $A_3$: es el factor lineal del polinomio de dimensión de Weyl de $A_3$, y falla ya
para $SU(5)$, donde el valor node-independiente es una suma de Littlewood--Richardson sin tal
factorización (por ejemplo toma el valor primo $11$ en los labels $(1,1,0,0)$, que ningún
producto de factores lineales reproduce a lo largo del retículo). Verificamos el
Lema~\ref{lem:node} en los tipos $A_2\!-\!A_4,\,D_4$ (node-independientes) y
$B_2,B_3,C_3,G_2,F_4$ (dos valores, divididos exactamente por raíz larga/corta, p.ej.\ $F_4$
dando $[21,21,15,15]$). El lema es elemental y probablemente folklore; lo incluimos porque
identifica precisamente cuál de las coincidencias de $SU(4)$ es genérica (la
node-independencia) y cuál excepcional (el producto), y esta última es uno de varios sentidos
en que $SU(4)$ es especial --- el tema de la siguiente sección.

% ===================================================================
\section{¿Por qué $SU(4)$? El rango crítico}\label{sec:rank}
¿Es el criterio un enunciado sobre el Modelo Estándar, o sobre $SU(4)$? Agranda el grupo, y la
respuesta separa ambos. El Teorema~\ref{thm:cell} es sobre $SU(4)$; si la misma rigidez
sobrevive para grupos mayores afila lo que ``$SU(4)$'' está haciendo en realidad. La celda
impone tres restricciones de carga, $Y(Q)=\tfrac16$, $Y(u)=\tfrac23$, $Y(d)=-\tfrac13$, sobre
un funcional de hipercarga que, por ser sin traza y $SU(2)_L$-invariante, tiene
$\mathrm{rango}(G)-1$ parámetros libres. El recuento es toda la historia:
\begin{center}\small
\begin{tabular}{lccl}
\toprule
\rowcolor{hdrblue}
\textcolor{white}{grupo} & \textcolor{white}{rango} & \textcolor{white}{parámetros de $Y$} &
\textcolor{white}{la celda}\\
\midrule
\rowcolor{softgrey}
$\boldsymbol{SU(4)}$ & $\mathbf 3$ & $\mathbf 2$ &
\textcolor{BrickRed}{\textbf{sobredeterminada}} --- condición aritmética rígida\\
$SU(5)$ & $4$ & $3$ & \textcolor{okgreen}{\textbf{marginalmente resoluble}} --- se disuelve\\
\rowcolor{softgrey}
$SU(6)$ & $5$ & $4$ & \textcolor{RoyalBlue}{\textbf{subdeterminada}} --- reps genéricas admiten\\
\bottomrule
\end{tabular}
\end{center}
Con tres restricciones y $\mathrm{rango}(G)-1$ parámetros, la celda es una obstrucción genuina
(un alineamiento de medida nula) exactamente cuando $\mathrm{rango}(G)-1<3$, i.e.\ en
$\mathrm{rango}=3$. Entre los grupos simples suficientemente grandes para contener la celda,
$SU(4)$ es el único caso marginal; en rango $4$ y superior una representación libre de
anomalías \emph{genérica} admite, siendo la única obstrucción residual el tener demasiado pocos
modos supervivientes. Lo hemos comprobado directamente: con una hipercarga libre, el $73\%$ de
las representaciones de $SU(5)$ y el $85\%$ de las de $SU(6)$ (hasta dimensión $200$) admiten,
en todas las clases del centro, mientras que $SU(4)$ permanece escaso y selectivo en el centro.
La aritmética ``rango $\ge$ algo'' es clásica; lo que la celda añade es la lectura de que
$SU(4)$ se sitúa exactamente en su margen.

\begin{figure}[t]
\centering
\includegraphics[width=0.68\linewidth]{fig_cell_dissolution_es.pdf}
\caption{La condición rígida de la celda se disuelve sobre el rango crítico. Con un funcional
de hipercarga libre ($\mathrm{rango}(G)-1$ parámetros contra tres restricciones), $SU(4)$
(rango $3$) está sobredeterminada y admite solo un conjunto escaso y selectivo en el centro;
$SU(5)$ y $SU(6)$ (rango $\ge4$) pasan a ser genéricamente resolubles, y una gran fracción de
representaciones admite en \emph{todas} las clases del centro (medido hasta dimensión $200$). La
transición es el recuento de \S\ref{sec:rank}, no una propiedad de un grupo concreto.}
\label{fig:dissolution}
\end{figure}

Hay una advertencia honesta que es también la física útil. ``Se disuelve en rango mayor''
supone que la hipercarga es \emph{libre} de ser cualquier combinación superviviente. En un
modelo real la hipercarga a menudo está \emph{fijada} --- fijada, digamos, por el requisito
$\sin^2\theta_W=\tfrac14$ que motiva estas construcciones de $SU(4)$. Fijar consume los
parámetros libres, y entonces la obstrucción puede volver a cualquier rango. Esto es justo lo
que sucede en el reciente modelo grand gauge--Higgs de $SU(7)$ de Komori y Maru, donde el
generador de hipercarga bulk, fijado para reproducir $\sin^2\theta_W$, asigna solo cargas
enteras y semienteras, ``de modo que las hipercargas fraccionarias de los quarks no pueden
obtenerse de los campos [bulk]'' y los quarks son forzados a una brana~\cite{KomoriMaru}.
Nuestro recuento lee su hecho entero/semientero como un resultado de
grados-de-libertad-consumidos: con la hipercarga fijada, no queda libertad de Cartan residual
para colocar la celda fraccionaria sobre tripletes de color bulk. La mitad robusta de la
historia es la que no depende de la suposición: $SU(4)$ es rígido \emph{incluso bajo máxima
libertad de hipercarga}, y por eso su criterio de celda es tan afilado como el
Teorema~\ref{thm:cell}.

\emph{Establecido aquí:} $SU(4)$ es el único rango en que la celda es una obstrucción rígida
bajo hipercarga libre, y el mismo recuento explica la obstrucción bulk-quark de rango mayor
como un efecto de libertad-consumida. \emph{Abierto:} la dependencia precisa del orbifold
$\ZZ_N$ y del $n$ general, expuesta a continuación.

% ===================================================================
\section{Reflexiones y propuestas}\label{sec:reflections}
Tres hilos parecen dignos de tirar.

\paragraph{La celda es el invariante.} A lo largo de los grupos anteriores, el álgebra gauge, la
representación, la dimensión y el centro cambian todos; la celda $\pm\tfrac12$ no. Es el
esqueleto de isospín débil de una generación, el mismo objeto en todo embedding, y el criterio
se lee mejor no como ``qué representaciones de $SU(4)$'' sino como ``¿contiene el grupo la celda
en el sector de centro correcto, con sitio para cerrarla?''. El centro suministra la aritmética,
la geometría suministra el sitio; la celda suministra el objetivo.

\paragraph{La dimensión es una segunda perilla.} El recuento factoriza como (ramificación de
Levi) $\times$ (proyección de orbifold). El primer factor es independiente de la dimensión; el
segundo es donde entra la compactificación. Una única $\ZZ_2$ actuando sobre el eje de
quiralidad parte los modos por la mitad; una $\ZZ_2$ bloqueada a un eje existente --- como lo
está la segunda paridad de $T^2/\mathbb{Z}_2$ sobre el sector de singletes, de modo que los
recuentos en cinco y seis dimensiones aquí coinciden --- no; un orbifold $\ZZ_N$ retiene una
fracción $1/N$. Proyecciones de dimensión mayor o de $N$ mayor por tanto \emph{encogen} el
recuento de modos, de modo que la obstrucción de pequeñez-geométrica (la única que queda en
rango gauge alto) puede revivir. El rango gauge y la proyección de orbifold son perillas
opuestas sobre la misma celda: una disuelve la condición aritmética, la otra aprieta la
geométrica. Hacer de la dependencia $\ZZ_N$ un teorema de recuento preciso es un paso siguiente
limpio.

\paragraph{Una forma cerrada más allá de $SU(4)$.} El Lema~\ref{lem:node} deja abierta la
dimensión de invariantes node-independiente exacta para $SU(n)$: es una suma de
Littlewood--Richardson cuya especialización a $SU(4)$ resulta ser un producto. Si admite una
forma cerrada de hook-content o determinantal --- y si los valores ``triangulares'' de $SU(5)$
$\binom{k+2}{2}=\dim\mathrm{Sym}^k(\mathbb C^{3})$ generalizan a
$\dim\mathrm{Sym}^k(\mathbb C^{n-2})$ para todas las potencias simétricas --- es una pregunta
combinatoria bien planteada que nuestros datos ya restringen.

% ===================================================================
\section{Conclusión}\label{sec:conclusion}
Hemos convertido la minimalidad observada en la Parte~I en una regla exacta. Una representación
$(a,b,c)$ de $SU(4)$ contiene una celda de quarks del Modelo Estándar sobre $T^2/\mathbb{Z}_2$
si y solo si su carga del centro es impar, su nodo central está excitado, y sus labels suman al
menos tres; la derivación corre por una ramificación de Levi cuya torre de Clebsch--Gordan,
proyectada sobre los modos cero quirales del orbifold, da el recuento cerrado
$N=(b+1)(a+c+1)/2$. Hemos sido deliberados con el crédito: la aritmética del centro, la
ramificación, y el recuento de parámetros de hipercarga son clásicos, y los citamos como tales;
lo nuestro es la proyección quiral que produce el recuento cerrado, el empaquetado en un único
criterio, y la lectura de la celda $\pm\tfrac12$ como el invariante que sitúa a $SU(4)$ en el
rango crítico. Lo último también explica, en vez de meramente reformular, por qué las
construcciones de rango mayor con hipercarga fijada no pueden mantener sus quarks en el bulk. El
criterio es exacto y verificado por máquina; las preguntas que abre --- la forma cerrada para
$n$ general y la dependencia dimensional $\ZZ_N$ --- están enunciadas con precisión suficiente
para ser zanjadas. Y cierra de vuelta sobre el hilo que la Parte~I dejó colgando: la viabilidad
se decide modelo a modelo, pero un modelo debe primero ser \emph{elegible}, y el
Teorema~\ref{thm:cell} es el test de elegibilidad --- dice qué representaciones pueden siquiera
entrar en la búsqueda de una generación viable de $SU(4)$ en seis dimensiones. De un testigo al
paisaje en que vive, ese es el arco de los dos trabajos.

\paragraph{Reproducibilidad.} Todo recuento se regenera desde los scripts anexos de SageMath:
la construcción directa de la celda~\eqref{eq:det}, la derivación por ramificación de Levi
de~\eqref{eq:count}, la comprobación exhaustiva del criterio hasta dimensión $900$, los barridos
de rango de $SU(5)/SU(6)$, y el barrido de node-independencia a lo largo de los tipos de Lie del
Lema~\ref{lem:node}.

\begin{thebibliography}{99}
\bibitem{PartI} C.~Mar\'in, \emph{Anomaly- and Tadpole-Compatible Fermion Completion of
Six-Dimensional $SU(4)$ Gauge--Higgs Unification} (Parte~I), Zenodo (2026),
DOI:\,10.5281/zenodo.21432625.
\bibitem{SSSW} C.A.~Scrucca, M.~Serone, L.~Silvestrini, A.~Wulzer, \emph{Gauge-Higgs
Unification in Orbifold Models}, JHEP \textbf{02} (2004) 049, hep-th/0312267.
\bibitem{AHMN} K.~Akamatsu, T.~Hirose, N.~Maru, A.~Nago, \emph{Two Higgs Doublet Model from
Gauge--Higgs Unification in Six Dimensions}, arXiv:2603.05857; y
\emph{Gauge--Higgs Unification of $SU(4)$ Model}, arXiv:2312.08608.
\bibitem{PatiSalam} J.C.~Pati, A.~Salam, \emph{Lepton Number as the Fourth Color},
Phys.\ Rev.\ D \textbf{10} (1974) 275; Erratum \emph{ibid.} \textbf{11} (1975) 703.
\bibitem{GeorgiGlashow} H.~Georgi, S.L.~Glashow, \emph{Unity of All Elementary Particle
Forces}, Phys.\ Rev.\ Lett.\ \textbf{32} (1974) 438.
\bibitem{SO10FM} H.~Fritzsch, P.~Minkowski, \emph{Unified Interactions of Leptons and
Hadrons}, Ann.\ Phys.\ \textbf{93} (1975) 193.
\bibitem{SO10G} H.~Georgi, \emph{The State of the Art --- Gauge Theories}, en
\emph{Particles and Fields 1974}, AIP Conf.\ Proc.\ \textbf{23} (1975) 575.
\bibitem{E6} F.~G\"ursey, P.~Ramond, P.~Sikivie, \emph{A Universal Gauge Theory Model Based
on $E_6$}, Phys.\ Lett.\ B \textbf{60} (1976) 177.
\bibitem{GeorgiFlavor} H.~Georgi, \emph{Towards a Grand Unified Theory of Flavor},
Nucl.\ Phys.\ B \textbf{156} (1979) 126.
\bibitem{Slansky} R.~Slansky, \emph{Group theory for unified model building},
Phys.\ Rept.\ \textbf{79} (1981) 1.
\bibitem{Yamatsu} N.~Yamatsu, \emph{Special Grand Unification}, PTEP \textbf{2017} 061B01,
arXiv:1704.08827.
\bibitem{HermsRuhdorfer} J.~Herms, M.~Ruhdorfer, \emph{How common are grand unified
theories?}, Phys.\ Rev.\ D \textbf{112} (2026) 115041, arXiv:2408.11089.
\bibitem{Hucks} J.~Hucks, \emph{Global structure of the standard model, anomalies, and
charge quantization}, Phys.\ Rev.\ D \textbf{43} (1991) 2709.
\bibitem{Saller} M.~Saller, \emph{The Central Correlations of Hypercharge, Isospin, Colour
and Chirality in the Standard Model}, Nuovo Cim.\ A \textbf{111} (1998) 1375,
arXiv:hep-th/9802112.
\bibitem{Tong} D.~Tong, \emph{Line Operators in the Standard Model}, JHEP \textbf{07}
(2017) 104, arXiv:1705.01853.
\bibitem{Alonso} R.~Alonso, D.~Dimakou, J.~West, \emph{Fractional-charge hadrons and
leptons to tell the Standard Model group apart}, arXiv:2404.03438.
\bibitem{ArandaWudka} A.~Aranda, J.~Wudka, \emph{Constraints on realistic gauge-Higgs
unified models}, Phys.\ Rev.\ D \textbf{82} (2010) 096005, arXiv:1008.3945.
\bibitem{KomoriMaru} K.~Komori, N.~Maru, \emph{$SU(7)$ Grand Gauge--Higgs Unification},
arXiv:2503.04090.
\bibitem{Fulton} W.~Fulton, \emph{Young Tableaux}, London Math.\ Soc.\ Student Texts
\textbf{35}, Cambridge University Press (1997).
\bibitem{GW} R.~Goodman, N.R.~Wallach, \emph{Symmetry, Representations, and Invariants},
Graduate Texts in Mathematics \textbf{255}, Springer (2009).
\bibitem{KawamuraGUT} Y.~Kawamura, \emph{Triplet-doublet splitting, proton stability and
extra dimension}, Prog.\ Theor.\ Phys.\ \textbf{105} (2001) 999, arXiv:hep-ph/0012125.
\end{thebibliography}

\end{document}
